% BJC Course Book preamble
% Adapted from snap-cloud/manual/_latex-template/template.tex
% (KOMA scrbook, US Letter, snap brand colors, snap heading style)

\documentclass[11pt,oneside]{scrbook}

\usepackage[paperheight=11in,
   paperwidth=8.5in,
   top=0.75in,
   bottom=0.75in,
   left=0.75in,
   right=0.75in,
   headsep=10pt,
   footskip=24pt,
   includehead, includefoot]{geometry}

\usepackage{iftex}
\ifPDFTeX
  \usepackage[T1]{fontenc}
  \usepackage[utf8]{inputenc}
  % Map a handful of common-in-curriculum unicode glyphs that pdflatex
  % can't render without explicit declarations. xelatex/lualatex
  % handle them natively.
  \DeclareUnicodeCharacter{2003}{\quad}      % em space
  \DeclareUnicodeCharacter{2002}{\enspace}   % en space
  \DeclareUnicodeCharacter{2009}{\,}         % thin space
  \DeclareUnicodeCharacter{200B}{}           % zero-width space
  \DeclareUnicodeCharacter{2192}{$\rightarrow$}
  \DeclareUnicodeCharacter{2190}{$\leftarrow$}
  \DeclareUnicodeCharacter{2191}{$\uparrow$}
  \DeclareUnicodeCharacter{2193}{$\downarrow$}
  \DeclareUnicodeCharacter{00F7}{$\div$}
  \DeclareUnicodeCharacter{00D7}{$\times$}
  \DeclareUnicodeCharacter{2212}{$-$}
  \DeclareUnicodeCharacter{2260}{$\neq$}
  \DeclareUnicodeCharacter{2264}{$\leq$}
  \DeclareUnicodeCharacter{2265}{$\geq$}
  \DeclareUnicodeCharacter{2026}{\ldots}
  \DeclareUnicodeCharacter{25A1}{\Box}
  \DeclareUnicodeCharacter{2610}{$\square$}
  \DeclareUnicodeCharacter{2611}{$\boxtimes$}
  \DeclareUnicodeCharacter{2153}{1/3}
  \DeclareUnicodeCharacter{2154}{2/3}
  \DeclareUnicodeCharacter{2155}{1/5}
  \DeclareUnicodeCharacter{2156}{2/5}
  \DeclareUnicodeCharacter{2157}{3/5}
  \DeclareUnicodeCharacter{2158}{4/5}
  \DeclareUnicodeCharacter{215B}{1/8}
  \DeclareUnicodeCharacter{215C}{3/8}
  \DeclareUnicodeCharacter{215D}{5/8}
  \DeclareUnicodeCharacter{215E}{7/8}
\fi
\ifXeTeX
  \usepackage{fontspec}
\fi
\ifLuaTeX
  \usepackage{fontspec}
\fi
\usepackage{lmodern}
\linespread{1.07}

\setkomafont{footnote}{\small}

\usepackage{graphicx}
\usepackage{grffile}
% Scale all images by 50% by default. The BJC source HTML routinely
% uses block screenshots and screen-resolution photos sized for a web
% browser, which print far too large in book form. Wrapping every
% \includegraphics in \scalebox{0.5} halves both dimensions regardless
% of which options pandoc emits.
\let\bjcorigincludegraphics\includegraphics
\renewcommand{\includegraphics}[2][]{%
  \scalebox{0.5}{\bjcorigincludegraphics[#1]{#2}}%
}
\usepackage{caption}
% Pandoc emits longtable + booktabs for any HTML <table>.
\usepackage{longtable}
\usepackage{booktabs}
\usepackage{array}
\usepackage{multirow}
% Pandoc-emitted helpers for table cells; some recent pandoc versions
% reference \real / \PandocStartIncludes which we provide as no-ops.
\providecommand{\real}[1]{#1}
\usepackage[dvipsnames,table]{xcolor}
\usepackage{changepage}
\usepackage{framed}
\usepackage[framemethod=tikz]{mdframed}
% amsmath is required for \text{} inside math mode (used by KaTeX
% expressions like \frac{a}{b} \text{ seconds}). Load before amssymb.
\usepackage{amsmath}
\usepackage{amssymb}
\usepackage[normalem]{ulem}
% Pandoc emits \ul for HTML <u> underline. Alias to ulem's \uline.
\let\ul\uline
% Some pandoc versions emit \st for HTML <s>/<strike>. ulem's \sout is
% the matching strikethrough.
\providecommand{\st}{\sout}
\usepackage{ifthen}
\usepackage{calc}
\usepackage{tikz}
\usepackage[absolute,overlay]{textpos}
\usepackage{enumitem}
\setlist[itemize]{noitemsep, topsep=0pt}
\setlist[enumerate]{topsep=0pt}

% Pandoc emits these for code highlighting / tightlist; provide harmless defaults.
\providecommand{\tightlist}{%
  \setlength{\itemsep}{0pt}\setlength{\parskip}{0pt}}
\providecommand{\passthrough}[1]{#1}

% Pandoc code-block scaffolding. pandoc wraps `<pre>` blocks in
% \begin{Shaded}\begin{Highlighting}[]...\end{Highlighting}\end{Shaded}
% along with a handful of "token-class" commands (\KeywordTok, etc.).
% We provide minimal definitions so the build doesn't fail when pages
% contain code samples; we render code in monospace without syntax color.
\usepackage{fancyvrb}
\usepackage{framed}
\definecolor{shadecolor}{RGB}{245,245,245}
\newenvironment{Shaded}{\begin{snugshade}}{\end{snugshade}}
\DefineVerbatimEnvironment{Highlighting}{Verbatim}{commandchars=\\\{\},fontsize=\small}
% pandoc highlighting token commands — all collapse to identity.
\newcommand{\VerbBar}{|}
\newcommand{\VERB}{\Verb}
\newcommand{\NormalTok}[1]{#1}
\newcommand{\KeywordTok}[1]{\textbf{#1}}
\newcommand{\DataTypeTok}[1]{#1}
\newcommand{\DecValTok}[1]{#1}
\newcommand{\BaseNTok}[1]{#1}
\newcommand{\FloatTok}[1]{#1}
\newcommand{\ConstantTok}[1]{#1}
\newcommand{\CharTok}[1]{#1}
\newcommand{\SpecialCharTok}[1]{#1}
\newcommand{\StringTok}[1]{\textit{#1}}
\newcommand{\VerbatimStringTok}[1]{#1}
\newcommand{\SpecialStringTok}[1]{#1}
\newcommand{\ImportTok}[1]{#1}
\newcommand{\CommentTok}[1]{\textit{\textcolor{black!60}{#1}}}
\newcommand{\DocumentationTok}[1]{\textit{#1}}
\newcommand{\AnnotationTok}[1]{\textit{#1}}
\newcommand{\CommentVarTok}[1]{\textit{#1}}
\newcommand{\OtherTok}[1]{#1}
\newcommand{\FunctionTok}[1]{#1}
\newcommand{\VariableTok}[1]{#1}
\newcommand{\ControlFlowTok}[1]{\textbf{#1}}
\newcommand{\OperatorTok}[1]{#1}
\newcommand{\BuiltInTok}[1]{#1}
\newcommand{\ExtensionTok}[1]{#1}
\newcommand{\PreprocessorTok}[1]{#1}
\newcommand{\AttributeTok}[1]{#1}
\newcommand{\RegionMarkerTok}[1]{#1}
\newcommand{\InformationTok}[1]{\textit{#1}}
\newcommand{\WarningTok}[1]{\textbf{#1}}
\newcommand{\AlertTok}[1]{\textbf{#1}}
\newcommand{\ErrorTok}[1]{\textbf{#1}}
\newcommand{\BuiltInModuleTok}[1]{#1}

% Snap brand colors (snap manual template)
\definecolor{snapblue}{RGB}{52,101,164}
\definecolor{snaporange}{RGB}{255,140,0}
\definecolor{snapcoverblue}{RGB}{52,101,164}

% BJC "div class" callout colors
\definecolor{learnbg}{RGB}{230,240,255}
\definecolor{takenotebg}{RGB}{255,245,220}
\definecolor{takenotebar}{RGB}{255,140,0}
\definecolor{forytdbg}{RGB}{225,245,225}
\definecolor{forytdbar}{RGB}{60,150,60}
\definecolor{endnotebg}{RGB}{245,245,245}
\definecolor{sidenotebg}{RGB}{250,250,235}
\definecolor{vocabbg}{RGB}{245,235,255}

% Heading style: snap manual rule beneath
\usepackage[explicit]{titlesec}
\newcommand{\snapheadingrule}{\vspace{-0.25\baselineskip}\noindent\rule{\linewidth}{0.4pt}}
\titleformat{\chapter}[hang]
  {\normalfont\LARGE\bfseries\color{snapblue}}{\thechapter\quad}{0pt}{#1}[\snapheadingrule]
\titlespacing*{\chapter}{0pt}{0pt}{12pt}
\titleformat{\section}[hang]
  {\normalfont\Large\bfseries}{\thesection\quad}{0pt}{#1}[\snapheadingrule]
\titleformat{\subsection}[hang]
  {\normalfont\large\bfseries\color{snapblue}}{\thesubsection\quad}{0pt}{#1}
\titleformat{\subsubsection}[hang]
  {\normalfont\normalsize\bfseries}{\thesubsubsection\quad}{0pt}{#1}
\titleformat{\paragraph}[runin]
  {\normalfont\normalsize\bfseries}{}{0pt}{#1}
\titleformat{\subparagraph}[runin]
  {\normalfont\normalsize\bfseries}{}{0pt}{#1}

\renewcommand*{\chaptermarkformat}{}

% Two-column TOC limited to chapter + section
\usepackage{multicol}
\setcounter{tocdepth}{1}
\setcounter{secnumdepth}{3}

% Index (vocab terms, On-the-AP-Exam codes, self-check questions).
% `imakeidx` lets us emit a two-column index in the TOC. The index is
% populated by \index{...} commands inserted by the LatexRenderer when
% it sees BJC's `.vocab` / `.exam` / `.assessment-data` markup.
\usepackage{imakeidx}
\makeindex[intoc, columns=2, title={Index}]

% Hyperref last, with linked TOC + bookmarks
\usepackage{hyperref}
\hypersetup{
  colorlinks,
  linkcolor={snapblue},
  citecolor={snapblue},
  urlcolor={snapblue},
  unicode=true,
  bookmarksnumbered=true,
  bookmarksopen=true,
}

% --- Callout environments mirroring BJC HTML classes --------------------
% mdframed (framemethod=tikz) is breakable across pages, so callouts
% containing tall images / many list items don't blow the LaTeX stack.
\newmdenv[
  linewidth=3pt,
  topline=false, bottomline=false, rightline=false,
  innerleftmargin=14pt, innerrightmargin=12pt,
  innertopmargin=10pt, innerbottommargin=10pt,
  skipabove=0.5\baselineskip, skipbelow=0.5\baselineskip,
  linecolor=snapblue, backgroundcolor=learnbg,
]{bjclearn}
\newmdenv[
  linewidth=3pt,
  topline=false, bottomline=false, rightline=false,
  innerleftmargin=14pt, innerrightmargin=12pt,
  innertopmargin=10pt, innerbottommargin=10pt,
  skipabove=0.5\baselineskip, skipbelow=0.5\baselineskip,
  linecolor=takenotebar, backgroundcolor=takenotebg,
]{bjctakenote}
\newmdenv[
  linewidth=3pt,
  topline=false, bottomline=false, rightline=false,
  innerleftmargin=14pt, innerrightmargin=12pt,
  innertopmargin=10pt, innerbottommargin=10pt,
  skipabove=0.5\baselineskip, skipbelow=0.5\baselineskip,
  linecolor=forytdbar, backgroundcolor=forytdbg,
]{bjcforyoutodo}
\newmdenv[
  linewidth=3pt,
  topline=false, bottomline=false, rightline=false,
  innerleftmargin=14pt, innerrightmargin=12pt,
  innertopmargin=10pt, innerbottommargin=10pt,
  skipabove=0.5\baselineskip, skipbelow=0.5\baselineskip,
  linecolor=black!30, backgroundcolor=endnotebg,
]{bjcendnote}
\newmdenv[
  linewidth=3pt,
  topline=false, bottomline=false, rightline=false,
  innerleftmargin=14pt, innerrightmargin=12pt,
  innertopmargin=10pt, innerbottommargin=10pt,
  skipabove=0.5\baselineskip, skipbelow=0.5\baselineskip,
  linecolor=takenotebar, backgroundcolor=sidenotebg,
]{bjcsidenote}
\newmdenv[
  linewidth=3pt,
  topline=false, bottomline=false, rightline=false,
  innerleftmargin=14pt, innerrightmargin=12pt,
  innertopmargin=10pt, innerbottommargin=10pt,
  skipabove=0.5\baselineskip, skipbelow=0.5\baselineskip,
  linecolor=snapblue, backgroundcolor=vocabbg,
]{bjcvocab}

% Callouts added for the wider class set (exam, atwork, dialogue,
% ifTime, takeItFurther, time, narrower variants).
\definecolor{exambg}{RGB}{255,240,225}
\definecolor{exambar}{RGB}{210,90,30}
\newmdenv[
  linewidth=3pt,
  topline=false, bottomline=false, rightline=false,
  innerleftmargin=14pt, innerrightmargin=12pt,
  innertopmargin=10pt, innerbottommargin=10pt,
  skipabove=0.5\baselineskip, skipbelow=0.5\baselineskip,
  linecolor=exambar, backgroundcolor=exambg,
]{bjcexam}

\definecolor{atworkbg}{RGB}{230,250,240}
\definecolor{atworkbar}{RGB}{40,130,90}
\newmdenv[
  linewidth=3pt,
  topline=false, bottomline=false, rightline=false,
  innerleftmargin=14pt, innerrightmargin=12pt,
  innertopmargin=10pt, innerbottommargin=10pt,
  skipabove=0.5\baselineskip, skipbelow=0.5\baselineskip,
  linecolor=atworkbar, backgroundcolor=atworkbg,
]{bjcatwork}

\definecolor{dialoguebg}{RGB}{240,240,250}
\definecolor{dialoguebar}{RGB}{90,90,180}
\newmdenv[
  linewidth=3pt,
  topline=false, bottomline=false, rightline=false,
  innerleftmargin=14pt, innerrightmargin=12pt,
  innertopmargin=10pt, innerbottommargin=10pt,
  skipabove=0.5\baselineskip, skipbelow=0.5\baselineskip,
  linecolor=dialoguebar, backgroundcolor=dialoguebg,
]{bjcdialogue}

% ifTime, takeItFurther, takeItTeased are shown unconditionally in the
% PDF (the web hides them behind <details>; the user wants them visible
% in the printed book). Distinct colors so readers can spot "optional"
% content at a glance.
\definecolor{iftimebg}{RGB}{225,235,255}
\definecolor{iftimebar}{RGB}{74,108,212}
\newmdenv[
  linewidth=3pt,
  topline=false, bottomline=false, rightline=false,
  innerleftmargin=14pt, innerrightmargin=12pt,
  innertopmargin=10pt, innerbottommargin=10pt,
  skipabove=0.5\baselineskip, skipbelow=0.5\baselineskip,
  linecolor=iftimebar, backgroundcolor=iftimebg,
]{bjciftime}

\definecolor{takefurtherbg}{RGB}{240,230,250}
\definecolor{takefurtherbar}{RGB}{143,86,227}
\newmdenv[
  linewidth=3pt,
  topline=false, bottomline=false, rightline=false,
  innerleftmargin=14pt, innerrightmargin=12pt,
  innertopmargin=10pt, innerbottommargin=10pt,
  skipabove=0.5\baselineskip, skipbelow=0.5\baselineskip,
  linecolor=takefurtherbar, backgroundcolor=takefurtherbg,
]{bjctakefurther}

\newenvironment{bjctime}{\begin{bjciftime}}{\end{bjciftime}}

% Narrower variants are layout helpers in the web view; render them as
% small page-break-safe indented blocks. mdframed with no borders gives
% us list-item-safe wrapping (adjustwidth + lists is fragile).
\newmdenv[
  linewidth=0pt,
  innerleftmargin=1em, innerrightmargin=1em,
  innertopmargin=2pt, innerbottommargin=2pt,
  skipabove=0.3\baselineskip, skipbelow=0.3\baselineskip,
  backgroundcolor=white,
]{bjcnarrower}
\newmdenv[
  linewidth=0pt,
  innerleftmargin=1em, innerrightmargin=1em,
  innertopmargin=2pt, innerbottommargin=2pt,
  skipabove=0.3\baselineskip, skipbelow=0.3\baselineskip,
  backgroundcolor=white, fontcolor=snapblue,
]{bjcnarrowblue}
\newmdenv[
  linewidth=0pt,
  innerleftmargin=1em, innerrightmargin=1em,
  innertopmargin=2pt, innerbottommargin=2pt,
  skipabove=0.3\baselineskip, skipbelow=0.3\baselineskip,
  backgroundcolor=white, fontcolor=takefurtherbar,
]{bjcnarrowpurple}

\newenvironment{bjcpseudop}{\par\medskip\noindent\ignorespaces}{\par\medskip}

% Inline icon: BJC has many small icons (load, partner). Render them at
% a fixed text-height so they don't blow up.
\newcommand{\bjcinlineicon}[1]{%
  \raisebox{-0.18\baselineskip}{\includegraphics[height=1.1\baselineskip]{#1}}%
}

% Snap! styling for the lightning bolt and run-link icon (best-effort,
% no fontawesome dependency here since we want to compile under pdflatex
% out of the box).
\newcommand{\snaplightning}{$\lightning$}
\providecommand{\lightning}{\ensuremath{\downarrow}}

% Snap! brand mark — matches the snap manual style (italic "!").
% HTML uses <span class="snap">snap</span> which CSS styles as "Snap!"
% with the trailing "!" in italic. We mirror that in print.
\newcommand{\snap}{Snap\textit{!}}

% Inline QR code for Snap! "run" links. \snapqr{path/to.png} renders
% the QR at ~3 baselines tall, raised so it sits next to the inline
% link text. The PNG itself is pre-generated by `qrencode` at build
% time and cached under out/qr/.
\newcommand{\snapqr}[1]{%
  \raisebox{-0.25\baselineskip}{%
    \bjcorigincludegraphics[height=2.6\baselineskip]{#1}%
  }%
}

% Image fallback box: when an image file can't be located we substitute
% a small placeholder so the build still succeeds.
\newcommand{\bjcmissingimg}[1]{%
  \fbox{\parbox{0.3\linewidth}{\centering\small\textit{[missing image:\\#1]}}}%
}

% Make a chapter that doesn't break to a new page (used for inter-unit
% labs/pages that are still numbered but shouldn't get a full page break
% per topic — keep default for now).
